\documentclass[xcolor=pdftex,svgnames,table]{beamer}
%\usepackage[table]{xcolor}
\usepackage{latexsym,amssymb,amsmath,stmaryrd}
\usepackage[utf8]{inputenc}
\usepackage[T1]{fontenc}
\usepackage[francais]{babel}
\usepackage{pgf}
\usepackage{tikz}
\usetikzlibrary{arrows,patterns,plotmarks,shapes,snakes,er,3d,automata,%
backgrounds,topaths,trees,petri,mindmap}
\usepackage{pgfbaseimage}
\usepackage{ifpdf}
\ifpdf
\usepackage{graphicx}
\else
%\usepackage[dvips]{graphicx}
\usepackage{pstricks,pst-tree,pst-node}
\newpsobject{showgrid}{psgrid}{subgriddiv=1,griddots=10,gridlabels=6pt}
\fi
\usepackage{moreverb}
\usepackage[normalem]{ulem}
\usepackage{version}
\usepackage{url}
\usepackage{multicol}
\usepackage{listings}
\lstset{
  language=[ANSI]C, 
%  gobble=2, 
%  escapeinside="", 
  basicstyle=\ttfamily, 
%  directivestyle=\color{Fuchsia},
  identifierstyle = \color{DarkOrange},
  keywordstyle=\color{DarkGreen}, 
  commentstyle=\color{red}, 
%  numbers=left, 
%  numbersep=5pt, 
%  numberstyle=\scriptsize, 
%  moredelim=[is][\color{gray}\itshape]{/*}{*/}, 
%  moredelim=[is][\alert]{/+}{+/}, 
%  morecomment=[is]{/=}{=/}, 
%  morecomment=[in]{comment=}{,} 
%  fancyvrb=true 
} 

\newcommand{\binaire}[1]{\ensuremath{\underline{#1}}}
\newcommand{\C}[1]{\texttt{#1}}
\newcommand{\bbbn}{\ensuremath{\mathbb{N}}}
%%%%%%%%%%%%%%%%%%%%%%%%%%%%%%%%%%%%%%%%%%%%%%%%%%%%%%%%%%%%%%%%
%% ccBeamer 0.1, 2007-07-02                                   %%
%% Written by Sebastian Pipping <webmaster@hartwork.org>      %%
%% ---------------------------------------------------------- %%
%% Licensed under Creative Commons Attribution-ShareAlike 3.0 %%
%% http://creativecommons.org/licenses/by-sa/3.0/             %%
%%%%%%%%%%%%%%%%%%%%%%%%%%%%%%%%%%%%%%%%%%%%%%%%%%%%%%%%%%%%%%%%


%% Images
\newcommand{\CcImageBy}[1]{%
	\includegraphics[scale=#1]{creative_commons/cc_by_30}%
}
\newcommand{\CcImageCc}[1]{%
	\includegraphics[scale=#1]{creative_commons/cc_cc_30}%
}
\newcommand{\CcImageDevNations}[1]{%
	\includegraphics[scale=#1]{creative_commons/cc_dev_nations_30}%
}
\newcommand{\CcImageNc}[1]{%
	\includegraphics[scale=#1]{creative_commons/cc_nc_30}%
}
\newcommand{\CcImageNd}[1]{%
	\includegraphics[scale=#1]{creative_commons/cc_nd_30}%
}
\newcommand{\CcImagePd}[1]{%
	\includegraphics[scale=#1]{creative_commons/cc_pd_30}%
}
\newcommand{\CcImageSa}[1]{%
	\includegraphics[scale=#1]{creative_commons/cc_sa_30}%
}
\newcommand{\CcImageSampling}[1]{%
	\includegraphics[scale=#1]{creative_commons/cc_sampling_30}%
}
\newcommand{\CcImageSamplingPlus}[1]{%
	\includegraphics[scale=#1]{creative_commons/cc_sampling_plus_30}%
}


%% Groups
\newcommand{\CcGroupBy}[1]{% zoom
	\CcImageBy{#1}%
}
\newcommand{\CcGroupByNc}[2]{% zoom, gap
	\CcImageBy{#1}\hspace*{#2}\CcImageNc{#1}%
}
\newcommand{\CcGroupByNcNd}[2]{% zoom, gap
	\CcImageBy{#1}\hspace*{#2}\CcImageNc{#1}\hspace*{#2}\CcImageNd{#1}%
}
\newcommand{\CcGroupByNcSa}[2]{% zoom, gap
	\CcImageBy{#1}\hspace*{#2}\CcImageNc{#1}\hspace*{#2}\CcImageSa{#1}%
}
\newcommand{\CcGroupByNd}[2]{% zoom, gap
	\CcImageBy{#1}\hspace*{#2}\CcImageNd{#1}%
}
\newcommand{\CcGroupBySa}[2]{% zoom, gap
	\CcImageBy{#1}\hspace*{#2}\CcImageSa{#1}%
}
\newcommand{\CcGroupDevNations}[1]{% zoom
	\CcImageDevNations{#1}%
}
\newcommand{\CcGroupNcSampling}[2]{% zoom, gap
	\CcImageNc{#1}\hspace*{#2}\CcImageSampling{#1}%
}
\newcommand{\CcGroupPd}[1]{% zoom
	\CcImagePd{#1}%
}
\newcommand{\CcGroupSampling}[1]{% zoom
	\CcImageSampling{#1}%
}
\newcommand{\CcGroupSamplingPlus}[1]{% zoom
	\CcImageSamplingPlus{#1}%
}


%% Text
\newcommand{\CcLongnameBy}{Attribution}
\newcommand{\CcLongnameByNc}{Attribution-NonCommercial}
\newcommand{\CcLongnameByNcNd}{Attribution-NoDerivs}
\newcommand{\CcLongnameByNcSa}{Attribution-NonCommercial-ShareAlike}
\newcommand{\CcLongnameByNd}{Attribution-NoDerivs}
\newcommand{\CcLongnameBySa}{Attribution-ShareAlike}

\newcommand{\CcNote}[1]{% longname
	This work is licensed under the \textit{Creative Commons #1 3.0 License}.%
}

\usetheme{classic}
\newcommand{\nowrite}{\put(10,-4){\includegraphics[scale=.05]{creative_commons/nopencil}}}
\newcommand{\youwrite}{\put(10,-4){\includegraphics[scale=.05]{creative_commons/pencil}}}
\newcommand{\writethat}{\includegraphics[scale=.05]{creative_commons/pencil}}
\newcommand{\aemporter}{\put(10,-6){\includegraphics[scale=.05]{creative_commons/szymonraj_Shopping_bag}}}


%%% Titre -- Cours 2
\title{Algorithmique et programmation -- Cours 2. Compilation. Instruction de contrôle \C{for}.}
\author{Pierre Boudes}
\date{\today}

\begin{document}

%% Page de titre et licence CC.
\begin{frame}
	\titlepage
	\vfill
	\begin{center}
		\CcGroupByNcSa{0.83}{0.95ex}\\[2.5ex]
		{\tiny\CcNote{\CcLongnameByNcSa}}
		\vspace*{-2.5ex}
	\end{center}
\end{frame}


\begin{frame}
  \frametitle{Langages et programmes (reloaded)}

\begin{columns}
\column[T]{4.5cm}
\structure{Humains}\\
\alert{langage naturel}
\column[T]{3.5cm}
\alert{Langage de programmation}
\begin{tikzpicture}
\useasboundingbox (0,0) rectangle (2,4.3);
\begin{scope}[shift={(6.5,3)}]%computer 3
\draw[ultra thick, rounded corners=4pt,fill=blue!15] (0,0) rectangle (1,0.7);
\draw[ultra thick, pattern=crosshatch dots] (0,0) -- (-0.4,-0.4) -- (0.6,
-0.4) -- (1, 0) -- cycle;  
\end{scope}
\begin{scope}[shift={(5,2.5)}]%computer 1
\draw[ultra thick, rounded corners=4pt,fill=blue!10] (0,0) rectangle (1,0.7);
\draw[ultra thick, pattern=crosshatch dots] (0,0) -- (-0.4,-0.4) -- (0.6,
-0.4) -- (1, 0) -- cycle;  
\end{scope}
\begin{scope}[shift={(5.6,1.9)}]%computer 2
\draw[ultra thick, rounded corners=4pt,fill=blue!20] (0,0) rectangle (1,0.7);
\draw[ultra thick, pattern=crosshatch dots] (0,0) -- (-0.4,-0.4) -- (0.6,
-0.4) -- (1, 0) -- cycle;  
\end{scope}
\begin{scope}[shift={(-3,2.5)}]%hum 1
  \draw[ultra thick,fill=pink] (0,0.9) circle (0.3cm); \draw[draw,line
  width=1.1mm,fill=pink, rotate=45, rounded corners=4pt] (-0.7, -0.15)
  rectangle (0.7, 0.15) (-0.15, -0.7) rectangle (0.15, 0.7) ;
  \draw[draw=none,fill=blue, rotate=45, rounded corners=4pt] (-0.7,
  -0.15) rectangle (0.7, 0.15) (-0.15, -0.7) rectangle (0.15, 0.7) ;
\end{scope}
\begin{scope}[shift={(-4,1)}]%hum 2
  \draw[ultra thick,fill=LightSalmon] (0,0.9) circle (0.3cm); \draw[draw,line
  width=1.1mm,fill=pink, rotate=45, rounded corners=4pt] (-0.7, -0.15)
  rectangle (0.7, 0.15) (-0.15, -0.7) rectangle (0.15, 0.7) ;
  \draw[draw=none,fill=DarkTurquoise, rotate=45, rounded corners=4pt] (-0.7,
  -0.15) rectangle (0.7, 0.15) (-0.15, -0.7) rectangle (0.15, 0.7) ;
\end{scope}
\begin{scope}[shift={(-2,1)}]%hum 3
  \draw[ultra thick,fill=DarkRed] (0,0.9) circle (0.3cm); \draw[draw,line
  width=1.1mm,fill=pink, rotate=45, rounded corners=4pt] (-0.7, -0.15)
  rectangle (0.7, 0.15) (-0.15, -0.7) rectangle (0.15, 0.7) ;
  \draw[draw=none,fill=DarkGreen, rotate=45, rounded corners=4pt] (-0.7,
  -0.15) rectangle (0.7, 0.15) (-0.15, -0.7) rectangle (0.15, 0.7) ;
\end{scope}
% \only<3-6>{
% \draw[draw,line width=4mm,->,yellow](-2,3)--node [black] (4,3) {utilisent} (4,3);
% }
% \only<4-6>{%programment
% \draw[draw,line width=4mm,->,Goldenrod](-1.9,2.5)--node [black] (4,3) {programment} (4.1,2.5);
% }
% \only<5-6>{%binaire
% \draw [draw=none,text=blue] (4.7,3.5) node {01101} -- +(-0.2,-0.4)
% node {10100};
% \draw [draw=none,text=blue] (7,1.5) node {10101} -- +(-0.5,-0.4) node
% {00101};
% \draw [draw=none,text=red] (-2,3.5) node {!!?} -- +(-0.2,-0.4)
% node {??};
% }
% \only<6-6>{%asm
% \draw [draw=none,text=DarkGreen] (5.1,4.1) node {\C{movb 0x65, \%eax}};
% \draw [draw=none,text=DarkGreen] (4.1,1.2) node {\C{movl \%eax, \%edx}};
% }

\draw[draw,line width=4mm,->,Goldenrod](-1.9,2.5)--node [black]
(0.4,3) {programment} (0.3,2.5);
\node [rectangle, draw, fill=blue!20, 
text width=3em, text centered, minimum height=4.6em] at (1.2,2.5) {};
\node [rectangle, draw, fill=blue!20, 
    text width=3em, text centered, minimum height=4.6em] at (1.1,2.6)
    {};
\node[line width=0mm,Goldenrod,rotate=60] at (1.1,2.6) {\Large\textbf{source}};
\node at (1.1,1.9) {\C{prog.c}};
\begin{scope}[shift={(2.4,2.2)},black!30]%computer 4
\draw[ultra thick, rounded corners=4pt,fill=blue!10] (0,0) rectangle (1,0.7);
\draw[ultra thick, pattern=crosshatch dots] (0,0) -- (-0.4,-0.4) -- (0.6,
-0.4) -- (1, 0) -- cycle;  
\end{scope}
\draw[draw,line width=4mm,->,blue!20](2,2.5)--node [black] (4,3)
{traduction} (4.1,2.5);
\only<2->{%algorithme
\node [draw, text badly ragged, cloud,cloud ignores aspect, cloud
puffs=11, fill=green!20,text width=2cm]
at (-3.1,3) {algorithmes\\ \alert{(pseudocode)}};
}
\end{tikzpicture}
\column[T]{4.5cm}
\structure{Ordinateurs}\\
\alert{langage machine}
\end{columns}
Pour le moment nous nous concentrons sur les rudiments de langage C.
\uncover<2->{
Mais ne perdons pas de vue que si savoir s'exprimer est nécessaire, il
faut aussi savoir réfléchir à ce que nous faisons faire à nos
programmes. \alert{Écrivez vos algorithmes !}
}
\end{frame}
 
%PLAN
\frame[label=plan]{\tableofcontents}


\section{Compilation}
\begin{frame}
  \frametitle{Compilation\nowrite}
\begin{itemize} 
\item  \structure{Compiler} un programme c'est traduire un texte (\structure{code
  source}) d'un langage de haut niveau (langage C) en code de bas
  niveau (\structure{code machine}), de manière à ce que le système
  d'exploitation puisse, au besoin, déclencher l'exécution de ce
  programme.\pause \textbf{Dans un langage compilé, l'étape de traduction a
    lieu une fois pour toutes.}\pause

\item \structure{Interpréter}, c'est faire en même temps la traduction et l'exécution du
texte d'un langage de haut niveau (un \structure{script}). Un interprète simule ainsi un
processeur capable d'exécuter le langage de haut niveau. \pause
\textbf{Dans un
langage interprété, l'étape de traduction a lieu à chaque exécution.}\pause
\item Le langage C est traditionnellement un langage compilé.
\end{itemize}
\end{frame}

\begin{frame}
  \frametitle{Compilation\nowrite}
 
Les cinq grandes étapes de la compilation :

  \begin{enumerate}
  \item Analyse lexicale
  \item Analyse syntaxique
  \item Analyse sémantique
  \item Génération du code
  \item Édition de liens
\end{enumerate}
\end{frame}

\subsection{Analyse lexicale}
\begin{frame}
  \frametitle{Analyse lexicale}
\begin{block}{Analyse lexicale}
 Identifie les \emph{lexèmes} (unités lexicales du langage). Les espaces sont inutiles (\C{3*x+1} ou \C{3 * x + 1}), sauf comme séparateurs (\C{int x}, \C{intx}).
\end{block}
\pause
Erreur lexicale :
\begin{description}
\item[code source] \C{int x = @;}
\item[compilation] \C{error: stray ‘@’ in program}
\end{description}
\pause
Erreur détectée uniquement au moment de l'analyse sémantique :\\
\begin{description}
\item[code source] \C{intx = 0;}
\item[compilation] \C{error: ‘intx’ undeclared (first use in this function)} 
\end{description}
\end{frame}

\subsection{Analyse syntaxique}
\begin{frame}
  \frametitle{Analyse syntaxique}
\begin{block}{Analyse syntaxique\youwrite} 
trouve la structure syntaxique, (arbre
  syntaxique), et teste l'appartenance au langage. 
\end{block}
\pause
Exemple : dans l'expression \C{x = 3 * x + 1}, est-ce que la
sous-suite  \C{x + 1} correspond à une structure syntaxique ? 
\pause
\begin{center}
\C{
\begin{tikzpicture} [level distance=5mm,every node/.style={fill=blue!20,circle,inner sep=1pt}]
\node {=}
child {node {x}} 
child {node {+}
  child {node {*}
    child {node {3}} 
    child {node {x}}
  }  
  child {node {1}} 
};
\end{tikzpicture}
}
\end{center}
\pause
\begin{description}
\item[code source] Un else sans if le précédant immédiatement (point-virgule mal placé ?)
\item[compilation] \C{error: expected expression before ‘else’} 
\end{description}

\end{frame}

\subsection{Analyse sémantique}
\begin{frame}
\frametitle{Analyse sémantique}
\begin{block}{Analyse sémantique\youwrite}
  trouver le sens des différentes actions
  voulues par le programmeur.
  \begin{itemize}
  \item Quels sont les objets manipulés par le programme,
  \item quelles sont les propriétés de ces objets,
  \item quelles sont les actions du programme sur ces objets.
  \end{itemize}

\end{block}
\pause
Beaucoup d'erreurs peuvent apparaître durant cette phase : identificateur utilisé mais non
déclaré (la réciproque génère un \emph{warning} avec l'option $-$\C{Wall}), opération n'ayant
aucun sens, etc.


\begin{description}
\item[code source] variable \C{x} utilisée mais non déclarée
\item[compilation] \C{error: ‘x’ undeclared (first use in this function)}
\end{description}
\end{frame}


\subsection{Génération du code}
\begin{frame}
\frametitle{Génération du code}
\begin{block}{Génération du code\youwrite}
  encodage en assembleur, optimisations et allocations des registres,
  traduction en \structure{code objet} (du code machine non exécutable
  en l'état). 
\end{block}
\end{frame}

\subsection{Édition de liens}
\begin{frame}
\frametitle{Édition de liens}
\begin{block}{Édition de liens\youwrite}
  le code objet des fonctions externes
  (bibliothèques) est ajouté à l'exécutable. Le point d'entrée dans le
  programme est choisi (\C{main}). Insertion de données de débogage (option \C{$-$g}).
\end{block}
\pause%\nowrite
\begin{description}
\item[code source] Oublie de \C{stdio.h} et utilisation de \C{printf}.
\item[compilation] \C{warning: incompatible implicit declaration of built-in function ‘printf’}
\end{description}
\pause
\begin{description}
\item[code source] Pas de fonction principale (\C{main})
\item[compilation] \C{Undefined symbols:  "\_main", ...}
\end{description}
\pause
À votre avis : \C{Undefined symbols: "\_prinft" } ?
\end{frame}


\section[gcc]{La compilation en pratique (gcc)}
\begin{frame}
  \frametitle{La compilation en pratique (gcc)}
\tikzstyle{block} = [rectangle, draw, fill=blue!20, 
    text width=7em, text centered, minimum height=2em,]
\tikzstyle{line} = [draw, -latex']
\tikzstyle{sline} = [draw, structure, dashed,  -latex']
\begin{tikzpicture}[node distance = 1.4cm, auto]
% les fichiers
%\matrix {
\node [block] (source) {\C{source.c}};
\node [right of = source] (nothing1) {};
\node [right of = nothing1] (nothing2) {};
\uncover<3->{\node [block, below of =nothing2] (pre) {\C{.c} (à l'écran)};
{\small \node [node distance=13em, draw, rounded corners,text
  width=10em, right of =pre] {Le \alert{préprocesseur} \hfill\youwrite\\ enlève les commentaires et exécute les instructions commençant par un dièse : \C{\#define} (rechercher-remplacer) et \C{\#include} (insertion de fichier).};}}
\uncover<4->{\node [block, below of =pre] (asm) {\C{source.i}};}
\uncover<5->{\node [block, below of =asm] (obj) {\C{source.o}};}
\node [block, below of =obj] (exe) {exécutable};
\node [node distance= .7cm, below of =exe] (gcc) {\C{gcc \uncover<2->{\alert{$-$Wall}} source.c $-$o nom} };
\uncover<7->{
\node [block, node distance = .1cm, above left of =source] (sourceb) {\C{source.c}};
\node [block, node distance = .1cm, above left of =pre] (preb) {\C{.c} (à l'écran)};
\node [block, node distance = .1cm, above left of =asm] (asmb) {\C{source.i}};
\node [block, node distance = .1cm, above left of =obj] (objb) {\C{source.o}};
}

% les arcs
\path [line] (source) |- (exe);
\uncover<3->{\path [line] (source) |- node {\C{gcc $-$E}} (pre);}
\uncover<4->{\path [line] (source) |- node {\C{gcc $-$S}} (asm);}
\uncover<5->{\path [line] (source) |- node {\C{gcc $-$c}} (obj);}
\uncover<3->{\path [sline] (source) -| node {\structure{Préprocesseur}} (pre);}
\uncover<4->{\path [sline] (pre) -- node {\structure{Compilation}} (asm);}
\uncover<5->{\path [sline] (asm) -- node {\structure{Génération du code objet}} (obj);}
\uncover<6->{\path [sline] (obj) -- node {\structure{Édition de liens}} (exe);}
\end{tikzpicture}
\end{frame}


\section[Le for du C]{L'instruction de contrôle for du langage C}
\subsection{Rappel sur la programmation structurée}
%%%%%%%%%%%%%%%%%%%%%


\begin{frame}
  \frametitle{Rappel sur la programmation structurée \nowrite}
    \begin{enumerate}
    \item Exécuter les blocs les uns à la suite des autres (\structure{séquence})
    \item si une certaine condition est vraie, exécuter un bloc sinon
      en exécuter un autre (\structure{sélection})
    \item recommencer l'exécution d'un bloc tant qu'une certaine
      condition est vraie (\structure{répétition}).
    \end{enumerate}
    Un bloc peut lui-même contenir une combinaison de blocs.
\pause
Aujourd'hui nous allons voir une première forme de répétition en C, le
\C{for} qui sert pour exprimer \youwrite:
\begin{itemize}
  \item la répétition un nombre fixé de fois \alert{répéter $n$ fois}
  \item l'itération sur un ensemble de cas (d'autres langages plus
    modernes ont le     \C{for each} pour itérer sur un ensemble
    d'éléments). \alert{Pour chaque entier de $1$ à $n$ faire \ldots}
\end{itemize}
\end{frame}

\subsection{Rappels sur l'instruction de contrôle if}
\begin{frame}[fragile]
  \frametitle{Rappels sur l'instruction de contrôle if \nowrite}

\structure{Syntaxe :} \C{if (\emph{condition}) \{\emph{\colorbox{green!20}{bloc1}}\} else 
\{\emph{\colorbox{red!20}{bloc2}}\}}.
\medskip

%\pause
  \begin{columns}
    \begin{column}[t]{5cm}
\structure{Code source}
\begin{verbatim}
/* avant */
if (age < 18)
{
  permis = 0;
}
else
{
  permis = 1;
}
/* après */
\end{verbatim}
   \end{column}
\pause
   \begin{column}[t]{4cm}
\structure{Schéma de traduction}\\[.3cm]
\tikzstyle{amil}=[rectangle, draw,  fill=blue!20, text badly ragged, text width=3.7cm]

\tikzstyle{line} = [draw, thick, -latex']
\begin{tikzpicture}[node distance=.7cm]
  \node [amil] (avant) {avant};
  \node [amil,below of=avant] (eval) {évalue la condition};
  \node [amil,below of =eval,fill=orange!20] (selection) {saute si elle est fausse};
 % \uncover<4->{%
    \node [amil,below of =selection,fill=green!20] (bloc1) {bloc1};
    \node [amil,below of =bloc1,fill=orange!20] (saut1) {saute après};
  %}
  \node [amil,below of =saut1,fill=red!20](bloc2) {bloc2};
  \node [amil,below of =bloc2] (apres) {après};
  \path [line] (selection) -| ++(2.6,0) |-  (bloc2);
%\uncover<4->{
\path [line] (saut1) -| ++(2.4,0) |-  (apres);
%}
\end{tikzpicture}
\end{column}
\end{columns}
\end{frame}

\subsection{L'instruction de contrôle for}
\begin{frame}[fragile]
  \frametitle{L'instruction de contrôle for\youwrite}

\structure{Syntaxe :}\\
 \C{for (\emph{instruct1}; \emph{condition}; \emph{\colorbox{yellow!20}{instruct2}}) 
\{\emph{\colorbox{green!20}{bloc}}\}}.
\medskip

\pause
  \begin{columns}
    \begin{column}[t]{6cm}
\structure{Code source}
\begin{verbatim}
/* avant */
for (i = 0; i < 5; i = i + 1)
{
  printf("%d\n", i);
  ...
}
/* après */
\end{verbatim}
La variable \C{i} est appelée \alert{variable de boucle}, elle doit être préalablement déclarée comme toute autre variable.

   \end{column}
\uncover<3->{
   \begin{column}[t]{4cm}
\structure{Schéma de traduction}\\[.3cm]
\tikzstyle{amil}=[rectangle, draw,  fill=blue!20, text badly ragged, text width=3.7cm]
\tikzstyle{line} = [draw, thick, -latex']
     \begin{tikzpicture}[node distance=.7cm]
       \node [amil] (avant) {avant};
       \node [amil,below of=avant] (instruction1) {instruction 1};
       \node [amil,below of =instruction1,fill=orange!20] (saut1) {saute à la condition};
       \node [amil,below of =saut1,fill=green!20] (bloc) {bloc};
\uncover<4->{\node [amil,below of =bloc,fill=yellow!20] (instruction2) {instruction 2};}
       \node [amil,below of =instruction2](eval) {évalue la condition};
       \node [amil,below of =eval,fill=orange!20](sautcond) {saute si elle est vraie};
       \node [amil,below of =sautcond] (apres) {après};
       \path [line] (saut1) -| ++(2.6,0) |-  (eval);
       \path [line] (sautcond) -| ++(2.4,0) |-  (bloc);
     \end{tikzpicture}
   \end{column}
}
 \end{columns}
\end{frame}

\subsection{Trace}
\begin{frame}[fragile]
  \frametitle{Trace\nowrite}
\begin{columns}
  \begin{column}[b]{4cm}
 \begin{lstlisting}[numbers=left,basicstyle=\ttfamily\scriptsize]
int main()
{
  /* Declaration et initialisation de variables */
 int i; /* var. de boucle */

  for (i = 0; i < 3; i = i + 1) /* pour chacune des 3 etapes */
  {
    printf("etape %d\n", i);
  }
  printf("i vaut %d\n", i);

  return EXIT_SUCCESS;  
}
\end{lstlisting}
\vspace{.4cm}
  \end{column}
\begin{column}[b]{4cm}
 \scriptsize
\pause
\rowcolors[\hline]{2}{green!25}{yellow!50} \arrayrulecolor{blue!75!gray}
  \begin{tabular}{|r|c|l|}
\hline
    ligne & \C{i} & sortie écran \\
\pause
   initialisation & ? &  \\
\pause
  6 & 0 & \\ 
\pause
  8 &  & \C{etape 0} \\
\pause
  9 & 1 & \\ 
\pause
  8 &  & \C{etape 1} \\
\pause
  9 & 2 & \\ 
\pause
  8 &  & \C{etape 2} \\
\pause
  9 & 3 & \\ 
\pause
  10 & & \C{i vaut 3} \\ 
\pause
12 & & Renvoie \C{EXIT\_SUCCESS} \\ 
\end{tabular}
  \end{column}
 \end{columns}
\end{frame}

\section{Algorithmes élémentaires}
\newcounter{savedenumi}

\begin{frame}
  \frametitle{Algorithmes : quels outils pour quels problèmes\nowrite}
\pause
\begin{columns}
 \begin{column}[t]{0.59\linewidth}
  \begin{enumerate}
    \item  Traiter des cas spécifiques 
      \begin{itemize}
      \item \alert{\C{if else}} (différencier)
      \item \alert{\C{\#define}} constantes symboliques (nommer)
 %     \item arbre de décision (organiser)
     \end{itemize}\pause
    \item Parcourir/générer des cas\pause
      \begin{itemize}
      \item \alert{boucle \C{for}} (rarement while)
%      \item tableaux
      \end{itemize}\pause
   \item Composer des cas
      \begin{itemize}
        \item boucles (parcourir/générer)
        \pause
      \item \alert{accumulateur} (à initialiser)
      \end{itemize}\pause
\setcounter{savedenumi}{\value{enumi}}
\end{enumerate}
\end{column}\hspace{-1cm}
\pause
\begin{column}[t]{0.59\linewidth}
\begin{enumerate}
\setcounter{enumi}{\value{savedenumi}}
   \item Sélectionner des cas
      \begin{itemize}
        \item boucles (parcourir/générer)
        \item \alert{\C{if}}  (sélectionner/traiter)
      \end{itemize}\pause
    \item[3'.] Dénombrer des cas
      \begin{itemize}
        \item boucles (parcourir/générer)
          \pause
        \item \alert{compteur} (à initialiser à 0)
      \end{itemize}
%     \item Rechercher un cas\pause
%       \begin{itemize}
%         \item boucle \alert{\C{while}}, conditions booléennes, \C{if}
%       \end{itemize}\pause
%     \item[2'.] Parcourir/générer : une \emph{ligne} mais aussi une \emph{surface}, un \emph{volume}...\pause
%       \begin{itemize}
%         \item imbriquer les boucles (var. $\neq$)
%       \end{itemize}\pause
%     \item Boucle événementielle
%   \begin{itemize}
% \item boucle \alert{while}
%   \end{itemize}\pause
%     \item Attente active 
%   \begin{itemize}
% \item boucle \alert{while}
%   \end{itemize}
\end{enumerate}
\end{column}
\end{columns}
\end{frame}

\section{Démos}
\begin{frame}
  \frametitle{Démos}
\end{frame}


\section{printf/scanf (1)}
\begin{frame}[fragile]
  \frametitle{printf/scanf (1) \nowrite}
\begin{itemize}
\item  Pour afficher un texte à l'écran, nous utilisons la fonction \alert{printf}
  (\emph{print formatted}). 
\item Chaque \verb+%+ 
dans le texte à afficher est substitué par la valeur formatée d'un
\alert{paramètre supplémentaire} de la fonction. Le caractère suivant le
symbole \verb+%+ 
signale le format à utiliser.  Un \verb+%d+
met une valeur au format \alert{entier décimal}.
\pause
\item Exemples :
  \begin{itemize}
\item
    \verb+printf("Bonjour\n")+ 
affiche \C{Bonjour} et un saut de
    ligne\pause
\item
    \verb+printf("i vaut %d\n", i)+ 
affiche \C{i vaut} suivi de la valeur décimale de i (et d'un saut de
ligne)\pause
\item
    \verb+printf("(%d, %d)\n", 31, -4)+ 
affiche \C{(31, -4)} et un saut de ligne.
 \end{itemize}
\pause
\item Réciproquement pour faire entrer dans le programme une donnée saisie par
  l'utilisateur, nous utiliserons \alert{\C{scanf}}.
\item Exemple : \verb+scanf("%d", &x)+
\end{itemize}
\end{frame}

\end{document}
